\documentclass{./worksheet}
\usepackage[colorlinks=true, urlcolor=blue]{hyperref}

\mytitle{Introduction to Javascript}

\begin{document}

This worksheet will give you some practice using the Javascript programming language. Make sure to refer to the slides if you are forgetting
any of the syntax or the library functions. You can also raise your hand and ask for help if you are stuck! \\

Javascript is a browser-based language, so you will be completing these tasks in
a playground like \url{https://runjs.app/play}. You can run JS locally on your computer
directly inside a browser or by installing a JS runtime like Node.js.
\begin{enumerate}

\item Write a Javascript arrow function that takes two integers as input and returns their difference.
Call that function on two integers of your choice and print a message to the screen,
in the format ``\verb|<int1>| minus \verb|<int2>| equals \verb|<diff>|".

\item \textbf{Challenge:} Create two variables and store a random integer between 0 and 100 in each one.
Use the \verb|Math.random()| function to generate these values. See \url{https://developer.mozilla.org/en-US/docs/Web/JavaScript/Reference/Global_Objects/Math/random}
for help.

\item Use the difference function on your two chosen or random integers to tell which of the two are bigger,
and print a message to the screen in the format ``\verb|<int1>| is bigger than \verb|<int2>|".

If they are the same, print a message in the format ``\verb|<int1>| equals \verb|<int2>|".

\item Modify the arrow function so that it returns an object with the two integers and their difference.

\item Create a new function that takes in a function as an argument, and:

1. Creates two random integers between 0 and 100 \\
2. Calls the function on the two integers \\
3. Returns the result

\item Call the new function passing in the difference function as its argument
(hint: you will need to store the result of the difference function in a variable first).

Print out the result in the format ``\verb|<int1>| minus \verb|<int2>| equals \verb|<diff>|".

Hint: use object destructuring to extract the values from the result object.

\item \textbf{Challenge:} Create an array of 5 strings in lowercase. Loop through the array
and print each string in uppercase.


\end{enumerate}


\end{document}
